\documentclass[11pt,letterpaper]{article}
\usepackage[margin=1in,headheight=15pt]{geometry}
\usepackage{fontspec}
\setmainfont{Latin Modern Roman}
\setsansfont{Latin Modern Sans}
\setmonofont[Scale=0.83]{DejaVu Sans Mono}
\usepackage{amsmath,amssymb,amsthm}
\usepackage{booktabs,longtable,array,enumitem}
\usepackage{fvextra}
\usepackage{xcolor}
\definecolor{ink}{HTML}{19384D}
\definecolor{muted}{HTML}{526575}
\usepackage[obeyspaces,spaces]{xurl}
\usepackage[colorlinks=true,linkcolor=ink,urlcolor=ink,citecolor=ink]{hyperref}
\usepackage{fancyhdr}
\pagestyle{fancy}
\fancyhf{}
\fancyhead[L]{\small\sffamily Further Lean proof simplifications}
\fancyhead[R]{\small\sffamily ProveIt / 26 October 2026}
\fancyfoot[C]{\thepage}
\setlength{\parindent}{0pt}
\setlength{\parskip}{0.52em}
\setlength{\emergencystretch}{2em}
\setlist{itemsep=0.22em,topsep=0.35em}
\fvset{fontsize=\small,breaklines=true,breakanywhere=true,tabsize=2,
  samepage=true,frame=single,rulecolor=\color{muted},framesep=6pt}
\urlstyle{tt}
\newcommand{\code}[1]{\nolinkurl{#1}}
\newcommand{\repo}[2]{\href{https://github.com/VladimirReshetnikov/ProveIt/blob/b310142330df89e48413c8fce62e7de223e1c7d5/#1}{#2}}
\newcommand{\modulelink}[2]{\repo{Computability/DiophantineConsistency/Lean/DiophantineConsistency/#1.lean}{#2}}
\newtheorem{proposition}{Proposition}[section]
\newtheorem{lemma}[proposition]{Lemma}
\title{\color{ink}\sffamily Further Lean Proof Simplifications\\[0.35em]
  \large Six additional ideas for \texttt{DiophantineConsistency}}
\author{Research and implementation report for Vladimir Reshetnikov}
\date{26 October 2026}

\begin{document}
\maketitle
\thispagestyle{empty}

\begin{abstract}
This report develops six concrete proof simplifications after auditing the five
archives in \code{docs/incoming} at the initial snapshot, the sixth retired
packet linked from its README, and two more packets added during this review.
The resulting changes replace two
unused internal inductions by an existing zero-or-successor decomposition;
reuse an exact Pell parity theorem; compose divisibility statements without
opening quotient witnesses; compare Pell indices by monotonicity in the base;
keep a nonsquare-norm argument entirely in the integers; and replace an explicit
finite-orbit construction by inverse invariance of a finite set.

The accompanying patch preserves the existing public theorem statements.
The report distinguishes local proof shortening, helper costs, dependency
simplification, and the scope of the Lean checks. The novelty claim is relative
to the inspected proposal packets. A promising evaluator refactor was excluded
because an earlier note already suggests a closely related semiring witness.
\end{abstract}

\vspace{0.4em}
\textbf{Pinned repository snapshot.}\quad
\href{https://github.com/VladimirReshetnikov/ProveIt/commit/b310142330df89e48413c8fce62e7de223e1c7d5}{\code{b310142330df89e48413c8fce62e7de223e1c7d5}}

\textbf{Verification versions.}\quad Lean 4.32.0; mathlib
\code{81a5d257c8e410db227a6665ed08f64fea08e997}.

\textbf{Closing audit snapshot.}\quad
\href{https://github.com/VladimirReshetnikov/ProveIt/commit/d163c476db9ca1142d2d73de8fa64c0b2c03a848}{\code{d163c476db9ca1142d2d73de8fa64c0b2c03a848}}.
The six target files and dependency pins are unchanged, and the patch also
passes its application check at this later snapshot.

\textbf{Recommended starting point.}\quad Apply the two unused-induction
deletions, the existing Pell parity theorem, and the integer norm argument
first. The Pell base comparison and finite-invariant-set argument are the most
useful mathematical changes to the proof structure.

\newpage
\setcounter{tocdepth}{1}
\tableofcontents
\newpage

\section{Outcome and measurement}

The deliverable is a reviewable six-file patch, complete replacement modules,
standalone comparison experiments, a prior-proposal exclusion ledger, and
validation evidence. The patch changes proofs and imports; it retains the
statements of the seven existing public theorems directly edited. The Pell
bootstrap file gains one reusable simultaneous monotonicity lemma.

\begin{longtable}{@{}p{0.07\textwidth}>{\raggedright\arraybackslash}p{0.32\textwidth}>{\raggedright\arraybackslash}p{0.55\textwidth}@{}}
\toprule
Idea & Main target & Actual simplification \\
\midrule
\endhead
1 & Object addition and multiplication at zero & Replace two specialized
\code{omega_ind} scripts and their private formula specifications by
\code{nat_zero_or_succ}. \\
2 & Pell auxiliary-index oddness & Reuse
\code{SignedPellResidue.odd_index_of_base_dvd} with the exact hypothesis
already present. \\
3 & Pell sign transport & Apply closure of divisibility under addition,
subtraction, and negation instead of constructing four quotient witnesses. \\
4 & Main versus first Pell index & Prove simultaneous monotonicity of the two
Pell coordinates in the base, then use strict monotonicity in the index. \\
5 & Nonsquare integer norm & Infer divisibility of square roots and cancel a
nonzero integer square; remove the detour through rational squares. \\
6 & Offset-19 exclusion & Use inverse invariance of a finite permutation's
forward-invariant even-state set in place of an explicit period and iteration
induction. \\
\bottomrule
\end{longtable}

\begin{center}
\begin{tabular}{@{}lrrr@{}}
\toprule
Module & Before & After & Removed \\
\midrule
\code{ObjectAdditionCancellation} & 118 & 96 & 22 \\
\code{ObjectNaturalZero} & 60 & 40 & 20 \\
\code{NativePellStepDown} & 248 & 220 & 28 \\
\code{NativePellBootstrap} & 685 & 677 & 8 \\
\code{MinimumNorm} & 308 & 304 & 4 \\
\code{NativeFieldCycles2011} & 111 & 103 & 8 \\
\midrule
\textbf{Total} & \textbf{1530} & \textbf{1440} & \textbf{90} \\
\bottomrule
\end{tabular}
\end{center}


The file totals in this table include comments, whitespace, helper lemmas,
and changed imports. They are reproducible source measurements, not a claim
about kernel term size or compilation speed. Separate proof-region counts in
the validation manifest explain where the savings occur. In particular, the
Pell monotonicity helper is charged to the change; the three-line consumer
proof is not advertised as a free replacement for the old proof.

The repository's requested Lean tree contains 574 Lean files at this snapshot.
The directory named \code{DiophantineConsistency} contains 391 of those files
and 71,601 physical lines, including several check modules. This review selects
a small number of proof changes from that larger codebase; it is not a
whole-project minimization result.

\section{How existing proposals were excluded}
\label{sec:novelty}

\subsection{The inspected material}

At the initial snapshot, the \repo{docs/incoming}{incoming directory} contained
these five archives:

\begin{enumerate}
\item \code{Lean_Proof_Simplification_2026-10-26.zip}
\item \code{ProveIt_Lean_Proof_Simplification_2026-10-26.zip}
\item \code{diophantine-lean-simplification-2026-10-26.zip}
\item \code{lean-proof-simplification-2026-10-26 (2).zip}
\item \code{lean-proof-simplification-2026-10-26.zip}
\end{enumerate}

Their 1,793 archive members contain no nested archives. The downloaded
archives' Git blob hashes match the repository directory entries. The audit
read their reports and relevant detailed notes and searched the proposed Lean
artifacts for overlap. Copies of original source files inside a packet were
treated as baseline evidence, rather than automatically treating every theorem
in such a copy as a proposal.

The incoming README also links a retired sixth packet,
\code{diophantine-lean-proof-simplification-2026-10-26.zip}, now placed under
\repo{Computability/DiophantineConsistency/research/lean-proof-simplification-20261026/PLACEMENT.md}{\code{research/lean-proof-simplification-20261026}}.
Its report and three detailed review notes were also inspected. Historical
entries concerning other research programs are outside this proof-simplification
comparison.

The bundled \code{audit/exclusion_ledger.json} groups 81 earlier idea families
and records 103 checked report or note locators. The readable ledger gives the
archive, member path, line, and a distinctive phrase for each entry. The
comparison is semantic: changing a proposed abstraction's name is not counted
as a new idea.

\subsection{Closing refresh after concurrent repository changes}

The repository advanced while this review was in progress. At the closing
snapshot \code{d163c476db9ca1142d2d73de8fa64c0b2c03a848}, the original five
archives had been placed in research directories, and two new archives were
present in \code{docs/incoming}:
\code{diophantine-lean-simplification-2026-10-26(1).zip} and
\code{lean-proof-simplification-review-2026-10-26.zip}. These were separately
audited before delivery. The closing addendum records their hashes, proposal
locations, and comparisons with the six retained ideas.

One new packet also proposes an edit in \code{NativePellStepDown}, but it
targets \code{eq_of_residues} and simplifies casts. That differs from the
odd-index theorem reuse and sign-transport changes delivered here. The audit
found no matching proposal for any of the six retained mechanisms in the
two new packets.

A Git comparison confirms that all six edited source files, the dependency
pins, and the imported repository source roots are unchanged between the
initial and closing snapshots. The patch passes a separate application check
against the later snapshot's index. The main Lean tree acquired additional
strip and machine modules; those additions are outside this patch and were
not claimed as freshly rebuilt endpoints.

\subsection{Topics deliberately left out}

Earlier packets already propose assignment naturality, finite witness layouts,
scanner iterations and fusion, guarded history folds, exact numerical Boolean
outputs, tag dispatch, sparse coefficient bounds, half-radix arithmetic,
finite-fiber sum identities, residual-square interfaces, semiring compiler
sharing, power-label certificates, indexed DAG equations, shared Split/Gap
circuit tails, source/window records, and transport of necessary definable
induction. They also contain smaller suggestions in ancillary notes. None of
those is offered here as an additional idea.

A concrete evaluator experiment was developed and then excluded. It constructs
the semiring evaluator's relational witness and invokes the existing
\code{evalRel_unique} theorem to simplify \code{ObjectNaturalSemiring.evalRel_iff}.
Although that exact consumer refactor is absent from the inspected packets,
\code{notes/syntax.txt:133} in the fifth archive already suggests a semiring
witness for evaluator totality. The strict novelty policy used for this
delivery excludes that neighboring development from the patch and from the
six contributions counted here.

\subsection{The novelty boundary of the six retained ideas}

The first idea eliminates two inductions whose step cases discard the
induction hypothesis. This differs from the earlier plan to transport a
necessary induction through a formula specification. The finite-permutation
idea uses inverse preservation of a finite state predicate, rather than the
earlier scanner-fold and iterate abstractions. The Pell comparison uses an
order theorem in the base parameter, rather than a uniform coefficient or
exponential envelope. The other three changes are exact local replacements
for which the audit found no corresponding proposal.

These are additional contributions relative to the eight inspected packets.
The report does not claim that basic mathematical facts such as integer
cancellation, monotonicity, or finite-set invariance are new mathematics.

\section{Idea 1: remove two specialized inductions}
\label{sec:no-induction}

\subsection{What the current proofs actually use}

\modulelink{ObjectAdditionCancellation}{\code{ObjectAdditionCancellation.lean}}
proves \code{addRel_zero_iff} by defining a private object-language formula,
proving its semantic specification, and applying \code{omega_ind}. The
successor case begins \code{intro n hn _}: its induction hypothesis is
discarded. The same pattern occurs in
\modulelink{ObjectNaturalZero}{\code{ObjectNaturalZero.mulRel_zero_iff}}.

The necessary elimination principle is already available:
\repo{SetTheory/BoundedConsistency/Lean/BoundedZFCConsistency/InternalSat.lean}{\code{nat_zero_or_succ}}
states that every member of a model's internal omega is either its zero or a
successor of another member. Importing its existing module and using this
decomposition removes the two private formulas and their specification proofs.

\subsection{Addition}

The relation \code{AddRel H x y (vempty H)} already supplies the naturalness
of its arguments. Apply \code{nat_zero_or_succ} to the second argument; no
new public hypothesis is needed.

If $y=0$, functionality of addition compares the given graph witness with the
zero law and yields $x=0$. If $y=\operatorname{succ}(n)$, totality supplies the
intermediate sum $b=x+n$. The successor law gives $x+y=\operatorname{succ}(b)$.
Functionality then forces the empty set to equal that successor, contradicting
the fact that $b$ is a member of its successor. No assertion about smaller
values of $y$ is used.

\VerbatimInput{snippets/addition-new.lean}

\subsection{Multiplication}

For multiplication, the zero case is immediate. In the successor case,
totality gives $b=xn$, and the successor equation identifies
$x\operatorname{succ}(n)$ with $b+x$. If that sum is zero, the additive zero
theorem yields $x=0$. Again, no induction hypothesis occurs in the argument.

\VerbatimInput{snippets/multiplication-new.lean}

\subsection{Semantic scope and exact savings}

The direct decomposition applies to all internal naturals, including
nonstandard elements of a model. There is no replacement of internal induction
by induction over the external type \code{Nat}. The existing decomposition
theorem itself was proved once by definable omega induction, so the improvement
removes two redundant specialized induction scripts rather than removing
induction from the foundational dependency chain.

The additive theorem retains its original unguarded statement. The
multiplicative theorem retains its two naturalness hypotheses. The changed
proof/helper blocks occupy 73 physical lines before the patch and 29 after it.
The explicit import of \code{BoundedZFCConsistency.InternalSat} in each file
makes the net full-file reduction 42 lines: $118\to96$ and $60\to40$.

\section{Idea 2: reuse the exact Pell parity theorem}

\subsection{The duplicated argument}

In \modulelink{NativePellStepDown}{\code{NativePellStepDown.index_residues}},
Pell classification yields an index $s$ with
\[
 Iu=X_I(s),\qquad I>1.
\]
The original proof then derives oddness of $s$ by a local even-index
congruence and divisibility argument. But
\modulelink{SignedPellResidue}{\code{SignedPellResidue.odd_index_of_base_dvd}}
already proves precisely that $I\mid X_I(s)$ forces $s$ to be odd when $I>1$.
The displayed equality supplies its hypothesis directly.

\begin{Verbatim}
  have hsodd : Odd s := SignedPellResidue.odd_index_of_base_dvd hI1 (by
    rw [← hIu]
    exact dvd_mul_right I u)
\end{Verbatim}

The new import is acyclic: \code{SignedPellResidue} depends on the common Pell
integer library, not on \code{NativePellStepDown}. The reduction also permits
removal of the unused local proof that $s>0$ and an unused name for the second
coordinate equality. Public parameters remain in place for compatibility.

\subsection{Why this is preferable to another abstraction}

No new parity interface or generic congruence framework is needed. This is
exact reuse of an existing theorem with the existing hypotheses. The
classification step, the odd-index polynomial identity, and the subsequent
step-down argument remain explicit. The change reduces the local proof's
arithmetic burden and the number of places that must be maintained if the Pell
congruence interface changes.

\section{Idea 3: keep divisibility witnesses abstract}

\subsection{The local sign problem}

Later in the same \code{index_residues} theorem, the proof knows
\[
 c\mid(s-p)\ \lor\ c\mid(s+p),
 \qquad J\equiv ts\pmod c,
 \qquad t=1\ \lor\ t=-1.
\]
It must derive $c\mid(J-p)$ or $c\mid(J+p)$. The original proof opens the
divisibility statements into integer quotient witnesses in each sign case,
constructs a new quotient, and checks its equality by
\code{linear_combination}.

Divisibility already has the required additive closure laws. Convert the
congruence once to $c\mid(ts-J)$ and compose it with the known divisibility
statement using \code{dvd_sub}, \code{dvd_add}, and \code{dvd_neg}. The branch
arithmetic is summarized below.

\begin{center}
\begin{tabular}{@{}cll@{}}
\toprule
$t$ & Divisible expressions used & Resulting divisible expression \\
\midrule
$1$ & $s-p$, $s-J$ & $(s-p)-(s-J)=J-p$ \\
$1$ & $s+p$, $s-J$ & $(s+p)-(s-J)=J+p$ \\
$-1$ & $s+p$, $-s-J$ & $-((s+p)+(-s-J))=J-p$ \\
$-1$ & $s-p$, $-s-J$ & $-((s-p)+(-s-J))=J+p$ \\
\bottomrule
\end{tabular}
\end{center}

The delivered proof uses a small local helper and \code{Or.imp} to map the two
branches. It does not introduce a new global ``congruence up to sign'' relation.
The delivered proof uses \code{convert!} so that the source expression is
elaborated without prematurely imposing the target expression. The resulting
ring identity then closes the conversion. An explicitly typed intermediate
divisibility statement is another readable way to control the same inference.

\VerbatimInput{snippets/sign-transport-new.lean}

\subsection{When this technique is useful}

Opening an existential divisibility witness is appropriate when the quotient
itself becomes later data. Here the quotient is immediately repackaged into
another proposition. Leaving it abstract makes the proof expose its actual
structure: an additive combination of elements of the principal ideal
$c\mathbb Z$. The argument retains both possible signs and adds no positivity
condition on the signed intermediate coordinate.

Ideas 2 and 3 affect separate regions of one theorem. Their net source savings
are reported together at the file level, so the two changes are not counted
twice.

The adjacent \code{signed_index_residues} wrapper receives the same treatment.
After the two cases of \code{Int.natAbs_eq}, it reuses the supplied congruence
\code{hVJ} or its negation directly. This reduces the absolute-value congruence
block from 11 to five lines. The identity for its square is obtained from
\code{Int.natCast_natAbs} and \code{sq_abs}; the argument never equates a
negative integer with its absolute value. These are supporting edits within
the sign-transport idea, not additional contributions.

\section{Idea 4: compare Pell bases instead of exponentials}
\label{sec:pell-monotone}

\subsection{The current comparison}

The original
\modulelink{NativePellBootstrap}{\code{NativePellBootstrap.main_index_gt_first}}
assumes $1<A$, $1<P$, and
\[
 c=Y_A(p),\qquad k=2Y_P(n),\qquad k<c.
\]
It also assumes $c>1$, $p>0$, $n>0$, and the relatively strong separation
$2A<2P-1$. Its 21-line proof first establishes lower bounds on the indices,
then compares
\[
 c\le(2A)^{p-1}\le(2A)^{n-1}
 <(2P-1)^{n-1}\le Y_P(n)\le k
\]
under the contradictory assumption $p\le n$.

The desired conclusion follows from a more direct order comparison.

\subsection{Simultaneous monotonicity}

For a natural parameter $a>1$, write $X_a(n)$ and $Y_a(n)$ for the two Pell
coordinates. Their first-order recurrences are
\begin{align*}
 X_a(0)&=1,&Y_a(0)&=0,\\
 X_a(n+1)&=aX_a(n)+(a^2-1)Y_a(n),
 &Y_a(n+1)&=X_a(n)+aY_a(n).
\end{align*}
All coefficients are nonnegative. If $1<a\le b$, simultaneous induction on
$n$ gives
\[
 X_a(n)\le X_b(n)\quad\text{and}\quad Y_a(n)\le Y_b(n).
\]
The simultaneous statement is the useful induction invariant: trying to
prove the inequality for the second coordinate alone would need information
about the first coordinate in the successor case.

\VerbatimInput{snippets/pell-monotonicity-new.lean}

The proof uses the positive first-order recurrence, avoiding the subtraction
that occurs in an alternative second-order Pell recurrence. Natural
subtraction in $a^2-1$ is handled by its ordinary monotonicity theorem.

\begin{proposition}
Let $1<A\le P$. If $Y_P(n)<Y_A(p)$, then $n<p$.
\end{proposition}
\begin{proof}
Base monotonicity gives $Y_A(p)\le Y_P(p)$. Hence
$Y_P(n)<Y_P(p)$. Strict monotonicity of the second coordinate in its index
then gives $n<p$.
\end{proof}

Since $Y_P(n)\le 2Y_P(n)=k<c=Y_A(p)$, the proposition proves the target.
The compatibility theorem keeps the old hypotheses but now has a three-line
body:

\VerbatimInput{snippets/pell-index-new.lean}

\subsection{A useful strengthening, with its cost made explicit}

The standalone experiment also checks the stronger theorem whose assumptions
are only $1<A$, $1<P$, $A\le P$, the two coordinate identities, and $k<c$.
The old hypotheses $c>1$, $p>0$, $n>0$, and strict base separation are not
needed for this logical implication. Existing callers retain their public
interface in the integration patch.

The old consumer proof has 21 body lines and the replacement has three. The
new monotonicity helper occupies 11 physical lines including its declaration
and documentation, and a separating blank line is retained. Thus this is a
small net reduction even when the reusable helper is charged immediately.
More significantly, it removes four exponential comparison steps, truncated
index subtraction, and two positivity subarguments from the consumer.

The current source still proves the stronger base-separation fact for its
other uses. Replacing that theorem, changing its public statement, or deleting
its proof is not part of this patch. A future cleanup can inspect its complete
caller set and decide whether the weaker base inequality is sufficient there.

\section{Idea 5: prove the norm fact entirely over the integers}

\subsection{The source argument and its replacement}

\modulelink{MinimumNorm}{\code{MinimumNorm.norm_ne_zero_of_not_isSquare}}
asserts
\[
 \neg\operatorname{IsSquare}(d),\quad y\ne0
 \quad\Longrightarrow\quad x^2-dy^2\ne0
 \qquad(d,x,y\in\mathbb Z).
\]
The original proof casts the equality into $\mathbb Q$, introduces $x/y$,
clears a denominator, and transfers rational squareness back to the integers.

An integral proof is direct. Suppose $x^2=dy^2=y^2d$. Then $y^2\mid x^2$.
The integer theorem \code{Int.pow_dvd_pow_iff} gives $y\mid x$, so write
$x=yk$. Substitution yields
\[
 y^2k^2=y^2d.
\]
Because $y\ne0$, its square is nonzero and can be cancelled. Thus $d=k^2$,
contradicting the hypothesis.

\VerbatimInput{snippets/norm-new.lean}

\subsection{Why the coefficient domain matters}

The decisive statement is the divisibility-of-powers theorem for integers.
This should not be generalized silently to an arbitrary integral domain:
in a domain that is not integrally closed, an element can acquire a square
root in its fraction field without having one in the domain. The delivered
proof keeps the original integer statement and its hypothesis $y\ne0$.

The full replacement file compiles after removing the direct imports of
\code{Mathlib.Data.Rat.Lemmas} and \code{Mathlib.Tactic.FieldSimp}. This is a
reduction in the module's direct dependency vocabulary. It does not prove that
no transitive dependency elsewhere in the import graph mentions rational
numbers, and no compilation-time improvement is claimed.

\section{Idea 6: use inverse invariance of a finite set}

\subsection{The finite-state proof obligation}

\modulelink{NativeFieldCycles2011}{\code{NativeFieldCycles2011.offset19_excludes_510}}
works with a permutation $\sigma$ of a finite type, a state value $S_i<512$,
and a bit $\epsilon_i\le1$, satisfying
\[
 S_{\sigma i}=(2S_i+\epsilon_i)\bmod512.
\]
The exceptional bit may equal one only at a single fixed source value
$s_\ast$. The existing predecessor argument shows that if some value were
$510$, then $s_\ast=511$. The existing local successor lemma then shows
\[
 S_i\equiv0\pmod2\quad\Longrightarrow\quad
 S_{\sigma i}\equiv0\pmod4.
\]

The original final step proves evenness along all forward iterates, obtains
a positive period, constructs its predecessor index $p-1$, and uses return
to the starting point to get the contradiction $510\equiv0\pmod4$.

\subsection{The shorter proof}

Let $E=\{i\mid S_i\equiv0\pmod2\}$. The successor lemma implies
$\sigma(E)\subseteq E$. The restriction of an injective map to the finite set
$E$ is surjective, so $\sigma^{-1}(E)\subseteq E$ as well. Consequently,
if $S_i=510$, the predecessor $\sigma^{-1}i$ is even. Apply the same successor
lemma to that predecessor. Its successor is $i$, so $S_i$ is divisible by
four, which contradicts $510\bmod4=2$.

This exact finite-permutation fact is already present in mathlib as
\href{https://github.com/leanprover-community/mathlib4/blob/81a5d257c8e410db227a6665ed08f64fea08e997/Mathlib/GroupTheory/Perm/Finite.lean}{\code{Equiv.Perm.perm_symm_on_of_perm_on_finite}}.
The new final part of the proof is:

\VerbatimInput{snippets/finite-invariant-new.lean}

The direct import changes from the periodic-point library to
\code{Mathlib.GroupTheory.Perm.Finite}. The local arithmetic lemmas and the
separate offset-20 theorem remain the same. The full-file reduction is eight
physical lines; the larger benefit is removal of the custom orbit induction,
period witness, and predecessor-index arithmetic.

\subsection{What cannot be dropped}

Finiteness is essential to this inference. On an infinite orbit indexed by
$\mathbb Z$, a bounded state can stay at $511$ for all negative times, take
the zero exceptional bit once, and then evolve through
\[
 510,508,504,496,480,448,384,256,0,0,\ldots.
\]
Use the exceptional bit one while remaining at $511$, and zero from the
transition to $510$ onward. This satisfies the recurrence and the fixed-source
gate with $s_\ast=511$ on the infinite shift permutation. Thus neither this
argument nor the original one licenses deleting the finiteness hypothesis.

The permitted finite $511$ self-loop also remains permitted. The theorem
excludes the value $510$, not all nonzero cycles.

\section{Verification and integration}
\label{sec:verification}

All six complete replacement modules compiled successfully together in one
combined module tree with Lean 4.32.0 and the exact pinned mathlib revision.
The unchanged repository dependencies were built from the pinned sources;
the exact-version mathlib cache was reused. The combined run loaded the new
addition proof when checking the multiplication file and the new step-down
module when checking the Pell bootstrap file.

The existing \code{NativeFieldCycles2011Checks} module also compiled against
the replacement. Its checks include the allowed nonzero $511$ self-loop,
the zero loop, use of the offset-19 exclusion, and the offset-20 statement
on an infinite type with an arbitrary map.

The selected axiom audit reports only
\code{propext}, \code{Classical.choice}, and \code{Quot.sound} for the seven
edited existing theorems. The new \code{coordinates_mono_base} lemma reports
only \code{propext}. No \code{sorryAx}, native-decision trust axiom, or new
axiom appears in those results. Separate experiments also checked the
stronger Pell comparison and the direct internal-natural proofs.

The combined patch passes \code{git apply --check} at the pinned commit.
The seven retained public declaration headers match the originals after
whitespace normalization. The compiler leaves only the expected warnings
about public parameters retained for compatibility: \code{hy} in the
step-down theorem and \code{hp}/\code{hn} in the index comparison.

These checks cover the six changed modules and the stated regression and
axiom checks. The entire ProveIt repository and every downstream endpoint
were not rebuilt. The two removed direct imports in \code{MinimumNorm} have
no proof-bearing direct consumer in this directory beyond the import and
axiom-audit aggregators inspected during the review.


\subsection{Reproducibility}

The patch is pinned to the repository commit on the title page. It is intended
to be applied at the repository root. The archive contains:

\begin{itemize}
\item the LaTeX source and this PDF;
\item one combined patch and complete replacement files;
\item independent Lean experiments for the new mathematical arguments;
\item a verification script using the repository's pinned Lean/Lake setup;
\item source hashes, source-line metrics, successful compiler logs, and axiom
  output;
\item the exclusion ledger, candidate novelty comparison, and independent
  review notes.
\end{itemize}

The verification script applies no patch itself. It checks the expected
replacement hashes in an already patched checkout, builds the six affected
modules, and runs the axiom audit. This keeps the review and mutation steps
explicit. A preparatory \code{git apply --check} command is documented in the
README before the ordinary \code{git apply} command.

\subsection{What the evidence supports}

The checked replacement modules demonstrate that their retained declarations
elaborate and their proofs are accepted with the pinned dependencies. The
standalone experiments make it easy to inspect the essential arguments without
the surrounding source file. The axiom report checks the selected results for
admitted proofs or extra logical assumptions.

Line counts are deterministic textual measurements. They do not measure the
size of a kernel proof term, the transitive size of imports, or the time saved
on a cold project build. Compilation times in local logs are diagnostic only:
the environment reused an exact-version package cache and was not a controlled
benchmarking setup.

\section{Further research questions}

\subsection{Where should Pell parameter monotonicity live?}

The patch places the new helper immediately above its first consumer for a
small, reviewable change. If other modules need it, the natural destination
is a lower Pell sequence interface, with public names for the two projections
of the simultaneous theorem. A follow-up can inventory genuine consumers
before adding more wrappers. It should compare the complete helper cost with
the proof reductions at those call sites.

\subsection{Can the base-separation obligations be weakened downstream?}

The strengthened standalone theorem needs only $A\le P$. Inspecting all uses
of \code{source_base_separation} may reveal callers whose strict separation
proofs are now unnecessary. This requires an actual dependency and caller
audit; this report does not claim that the existing separation theorem can
already be deleted. The present patch preserves its public interface.

\subsection{Which internal inductions have no inductive content?}

The two successful examples suggest a targeted audit of \code{omega_ind}
uses whose step binder is literally discarded. The test is semantic: a proof
may hide an induction hypothesis in automation or use it under a different
name. Where only zero/successor elimination is needed, the existing internal
decomposition theorem should replace the formula-specific induction. Necessary
internal inductions must retain their definability witnesses; unrestricted
induction over arbitrary external predicates would change the mathematics.

\subsection{Can signed divisibility endpoints share a small algebraic lemma?}

The delivered patch uses a local helper because there is one identified
consumer. If a search finds several independent endpoints with the same
$\pm$-transport pattern, a public theorem about divisibility up to sign could
make them shorter. Its statement should encode the orientation of
\code{Int.ModEq} explicitly and preserve all sign possibilities. A global API
is warranted only after that reuse is demonstrated.

\subsection{Can other finite-state proofs use predicate invariance?}

Some finite-permutation arguments need an actual period or orbit length; those
should keep the periodic-point interface. When the only use of periodicity is
to return a forward-invariant property to a predecessor, the finite-set lemma
is a better fit. The relevant search should inspect the proof's consumer of
the period witness, not replace every finite orbit argument mechanically.

\subsection{How should proof simplification be scored?}

A useful follow-up benchmark can record four separate quantities: edited
source lines including helpers, proof-term size, import footprint, and
elaboration time under controlled conditions. A proof can improve one without
improving all four. The present work establishes source and structural
simplifications, with formal checking of the edited modules, and leaves
performance claims to such a benchmark.

\appendix
\section{Source map}

All repository links below resolve to the pinned snapshot. The exact baseline
and replacement file hashes are in \code{validation/manifest.json}.

\begin{longtable}{@{}>{\raggedright\arraybackslash}p{0.36\textwidth}>{\raggedright\arraybackslash}p{0.58\textwidth}@{}}
\toprule
Source & Role in this report \\
\midrule
\endhead
\modulelink{ObjectAdditionCancellation}{\code{ObjectAdditionCancellation.lean}} &
Original zero-sum induction and retained public statement. \\
\modulelink{ObjectNaturalZero}{\code{ObjectNaturalZero.lean}} &
Original zero-product induction and its downstream instance. \\
\repo{SetTheory/BoundedConsistency/Lean/BoundedZFCConsistency/InternalSat.lean}{\code{InternalSat.lean}} &
Internal zero-or-successor decomposition; no standardness assumption. \\
\modulelink{SignedPellResidue}{\code{SignedPellResidue.lean}} &
Existing exact odd-index theorem. \\
\modulelink{NativePellStepDown}{\code{NativePellStepDown.lean}} &
Auxiliary parity and signed divisibility transport. \\
\modulelink{NativePellBootstrap}{\code{NativePellBootstrap.lean}} &
Original exponential comparison of the main and first indices. \\
\repo{Computability/HilbertTenthProblem/Lean/Diophantine/Common/Pell.lean}{\code{Diophantine/Common/Pell.lean}} &
Pell recurrence and monotonicity interface used by the new base comparison. \\
\modulelink{MinimumNorm}{\code{MinimumNorm.lean}} &
Original rational-square norm argument. \\
\modulelink{NativeFieldCycles2011}{\code{NativeFieldCycles2011.lean}} &
Existing predecessor and successor arithmetic for the offset exclusion. \\
\href{https://github.com/leanprover-community/mathlib4/blob/81a5d257c8e410db227a6665ed08f64fea08e997/Mathlib/GroupTheory/Perm/Finite.lean}{\code{Mathlib/GroupTheory/Perm/Finite.lean}} &
Finite forward-invariant sets are invariant under the inverse permutation. \\
\repo{docs/incoming/README.md}{\code{docs/incoming/README.md}} &
Current and retired proposal packet inventory. \\
\bottomrule
\end{longtable}

\section{Reading the audit artifacts}

The exclusion ledger is intentionally more detailed than the prose summary.
An entry records an earlier idea family, its targets, and one or more precise
locations in archived material. It is a reference for overlap screening, not
an additional list of proposed work. The novelty comparison marks the
canonical-evaluation candidate as excluded and the unused constructor-algebra
design as not delivered.

The independent review notes examine logical scope: preservation of
nonstandard internal naturals, dependence on finite permutations, retention of
the $511$ self-loop, sign orientations, and the absence of circular use of the
theorems being replaced. Compiler acceptance and this conceptual review serve
different purposes; both are included so the intended simplification is easy
to assess.

\end{document}
